% Homework 2 -- Lecture 2: conditional probability, independence and Bayes' theorem. Paper only.
% The data come from hw02_data.tex, written by hw02_numbers.py (every number computed and unit-tested).
\par
% generated by hw02_numbers.py -- do not edit
\newcommand{\hwMOneShare}{70}
\newcommand{\hwMOneDef}{2}
\newcommand{\hwMTwoShare}{30}
\newcommand{\hwMTwoDef}{5}
\newcommand{\hwItems}{1{,}000}
\newcommand{\hwRed}{5}
\newcommand{\hwBlue}{3}
\newcommand{\hwAttOnePass}{102}
\newcommand{\hwAttOneFail}{18}
\newcommand{\hwAttTwoPass}{44}
\newcommand{\hwAttTwoFail}{36}
\newcommand{\hwSesOnePass}{73}
\newcommand{\hwSesOneFail}{27}
\newcommand{\hwSesTwoPass}{73}
\newcommand{\hwSesTwoFail}{27}
\newcommand{\hwHAEasyN}{200}
\newcommand{\hwHAEasyS}{190}
\newcommand{\hwHAHardN}{50}
\newcommand{\hwHAHardS}{30}
\newcommand{\hwHBEasyN}{50}
\newcommand{\hwHBEasyS}{49}
\newcommand{\hwHBHardN}{200}
\newcommand{\hwHBHardS}{130}
\newcommand{\hwBaseN}{2}
\newcommand{\hwSens}{95}
\newcommand{\hwFalsePos}{3}
\newcommand{\hwPeople}{100{,}000}
\newcommand{\hwSpam}{40}
\newcommand{\hwFreeSpam}{30}
\newcommand{\hwFreeHam}{2}
\newcommand{\hwMatch}{1{,}000{,}000}
\newcommand{\hwAccuracy}{99}
\newcommand{\hwSockB}{6}
\newcommand{\hwSockW}{4}
\newcommand{\hwSockT}{10}
\newcommand{\hwErrTotalWrong}{0.70}
\newcommand{\hwAttOneRate}{0.85}
\newcommand{\hwAttTwoRate}{0.55}


\clearpage
\section*{Homework 2: Conditional Probability, Independence and Bayes}
\noindent\textbf{Lecture 2.} Work on paper, by hand: a pen, a ruler and a basic calculator only. Show every step.
Keep fractions exact until the last line. Draw the tree or the table whenever a problem has stages or two
attributes; most mistakes in conditional probability are a wrong denominator, and a picture shows the right one.

\subsection*{Task 1. Knowing something shrinks the sample space}
\begin{enumerate}[label=(\alph*)]
  \item A fair die is rolled. Find $P(\text{at least } 5)$. A friend peeks and says the result is odd; find
        $P(\text{at least } 5 \mid \text{odd})$ by listing the outcomes that are still possible.
  \item One card is drawn from a deck of 52. Find $P(\text{king} \mid \text{face card})$ and
        $P(\text{face card} \mid \text{king})$ (the face cards are J, Q, K). Why are the two numbers different?
  \item Two dice are rolled. Using the $6 \times 6$ grid, find $P(\text{sum} = 8 \mid \text{first die} = 3)$ and
        $P(\text{first die} = 3 \mid \text{sum} = 8)$. Circle the cells of the new sample space each time.
  \item With the same two dice, find $P(\text{a double} \mid \text{sum} = 8)$.
\end{enumerate}

\subsection*{Task 2. A tree diagram}
A factory has two machines. Machine 1 makes \hwMOneShare\,\% of all items, and \hwMOneDef\,\% of its items are
defective. Machine 2 makes \hwMTwoShare\,\% of all items, and \hwMTwoDef\,\% of its items are defective.
\begin{enumerate}[label=(\alph*)]
  \item Draw the tree: first the machine, then defective or not. Write the probability on every branch.
  \item Multiply along each path. Check that the four path probabilities add up to 1, and explain why they must.
  \item Turn the paths into counts for \hwItems{} items: how many come from each machine, and how many of those are
        defective?
\end{enumerate}

\subsection*{Task 3. Drawing without replacement}
A bag holds \hwRed{} red and \hwBlue{} blue balls. Two balls are drawn one after the other, without putting the first
one back.
\begin{enumerate}[label=(\alph*)]
  \item Draw the tree and find $P(\text{both red})$ and $P(\text{both the same colour})$.
  \item Find $P(\text{the second ball is red})$, using both branches of the first draw. Compare it with
        $P(\text{the first ball is red})$. Surprised? Explain why the answer makes sense without any calculation.
  \item Find $P(\text{both red})$ again, but now the first ball is put back before the second draw. Which answer is
        larger, and why?
\end{enumerate}

\subsection*{Task 4. Contingency tables and independence}
Two hundred students took the same exam. The first table records whether they attended the lectures regularly; the
second, whether they sat the exam in the morning or the afternoon session.
\begin{center}
\begin{tabular}{l|rr|r}
 & Passed & Failed & Total \\ \hline
Attended regularly & \hwAttOnePass & \hwAttOneFail & \\
Did not attend & \hwAttTwoPass & \hwAttTwoFail & \\ \hline
Total & & & 200
\end{tabular}

\medskip
\begin{tabular}{l|rr|r}
 & Passed & Failed & Total \\ \hline
Morning session & \hwSesOnePass & \hwSesOneFail & \\
Afternoon session & \hwSesTwoPass & \hwSesTwoFail & \\ \hline
Total & & & 200
\end{tabular}
\end{center}
\begin{enumerate}[label=(\alph*)]
  \item Fill in the margins of both tables.
  \item From the first table find the joint probability $P(\text{attended} \cap \text{passed})$, the marginal
        $P(\text{passed})$, and the conditionals $P(\text{passed} \mid \text{attended})$ and
        $P(\text{passed} \mid \text{did not attend})$. Say which denominator you used for each.
  \item Is passing independent of attending? Check it two ways: by comparing a conditional with the marginal, and by
        comparing each cell with the count $\frac{\text{row total} \times \text{column total}}{n}$ that independence
        predicts.
  \item Do the same check for the session table. What do you conclude?
  \item Find $P(\text{attended} \mid \text{passed})$ and explain in words how it differs from
        $P(\text{passed} \mid \text{attended})$.
\end{enumerate}

\subsection*{Task 5. Simpson's paradox: two hospitals}
Two hospitals report their operations of two kinds, easy and hard:
\begin{center}
\begin{tabular}{l|rr|rr}
 & \multicolumn{2}{c|}{Hospital A} & \multicolumn{2}{c}{Hospital B} \\
 & operations & survived & operations & survived \\ \hline
Easy & \hwHAEasyN & \hwHAEasyS & \hwHBEasyN & \hwHBEasyS \\
Hard & \hwHAHardN & \hwHAHardS & \hwHBHardN & \hwHBHardS
\end{tabular}
\end{center}
\begin{enumerate}[label=(\alph*)]
  \item Compute each hospital's survival rate for easy operations, for hard operations, and overall.
  \item Which hospital is better at easy operations? At hard ones? Overall? Describe the paradox in one sentence.
  \item Explain where the reversal comes from. What share of each hospital's operations are hard?
  \item You need a hard operation. Which hospital do you choose, and which numbers would you show a friend who only
        read the overall rates?
\end{enumerate}

\subsection*{Task 6. Total probability and Bayes' theorem}
\begin{enumerate}[label=(\alph*)]
  \item \emph{The factory again (Task 2).} Use the law of total probability to find $P(\text{defective})$. An item is
        found to be defective: use Bayes' theorem to find the probability that machine 2 made it. Machine 2 makes
        fewer items -- why is it still the more likely culprit?
  \item \emph{Screening.} A disease affects \hwBaseN{} people in 1{,}000. A test detects it in \hwSens\,\% of ill
        people, and wrongly gives a positive result for \hwFalsePos\,\% of healthy people.
        \begin{enumerate}[label=(\roman*)]
          \item Imagine \hwPeople{} people. Fill in a table of counts: ill or healthy, positive or negative.
          \item A person tests positive. Find $P(\text{ill} \mid \text{positive})$ from your counts, then check it with
                Bayes' formula. Which group produces most of the positive results?
          \item The person takes a second, independent test of the same kind, and it is positive too. Use your answer
                to (ii) as the new prior and find the new posterior.
        \end{enumerate}
  \item \emph{A spam filter.} \hwSpam\,\% of incoming emails are spam. The word ``free'' appears in \hwFreeSpam\,\% of
        spam emails and in \hwFreeHam\,\% of the others. Name the prior and the likelihoods, and find
        $P(\text{spam} \mid \text{``free''})$.
\end{enumerate}

\subsection*{Task 7. Find the errors}
Each statement contains one error. Find it, explain it, and correct it where a number is involved.
\begin{enumerate}[label=(\alph*)]
  \item A prosecutor: ``Only one innocent person in \hwMatch{} would match this DNA sample. So the chance that the
        accused is innocent is one in \hwMatch.''
  \item ``From the attendance table (Task 4): $P(\text{passed} \mid \text{attended}) = \frac{\hwAttOnePass}{\text{number
        who passed}}$.''
  \item ``A drawer has \hwSockB{} black and \hwSockW{} white socks. I take two in the dark, so
        $P(\text{both black}) = \frac{\hwSockB}{\hwSockT} \cdot \frac{\hwSockB}{\hwSockT}$.''
  \item ``This test is \hwAccuracy\,\% accurate. My result is positive, so I am ill with probability
        \hwAccuracy\,\%.''
  \item ``In the attendance table, the pass rate is $\hwAttOneRate$ among those who attended and $\hwAttTwoRate$ among the others, so the
        overall pass rate is $\frac{\hwAttOneRate + \hwAttTwoRate}{2} = \hwErrTotalWrong$.''
  \item ``$P(B \mid A) = P(B)$, so $A$ and $B$ can never happen together: they are mutually exclusive.''
\end{enumerate}

\subsection*{Task 8. Explain it in words}
Answer each in three or four sentences.
\begin{enumerate}[label=(\alph*)]
  \item $P(\text{wet grass} \mid \text{it rained})$ is close to 1. Is $P(\text{it rained} \mid \text{wet grass})$ close
        to 1 too? What else could make the grass wet, and how does that enter Bayes' theorem?
  \item Why is screening a mostly healthy population for a rare disease hard, even with a good test? Use your counts
        from Task 6(b).
  \item Give an everyday example where pooling two groups could reverse a comparison, as in Task 5.
  \item \emph{Monty Hall.} A prize is behind one of three doors. You pick door 1. The host, who knows where the prize
        is and always opens an empty door you did not pick, opens door 3. Should you switch to door 2? Explain with
        conditional probability (a tree helps).
  \item A spam filter assumes that, \emph{given} the class (spam or not), the words of an email are independent.
        Explain what that assumption means, and give a pair of words for which it is clearly false.
\end{enumerate}

\subsection*{Optional challenge}
Prove each statement using only the definition $P(B \mid A) = \frac{P(A \cap B)}{P(A)}$ for $P(A) > 0$ and the rules
of Lecture 1.
\begin{enumerate}[label=(\alph*)]
  \item For a fixed event $A$ with $P(A) > 0$, $P(\,\cdot \mid A)$ is itself a probability: $P(B \mid A) \ge 0$,
        $P(A \mid A) = 1$, and $P(B \cup C \mid A) = P(B \mid A) + P(C \mid A)$ when $B \cap C = \emptyset$.
  \item Bayes' theorem $P(A \mid B) = \frac{P(B \mid A)\,P(A)}{P(B)}$.
  \item The law of total probability: if $A_1, \dots, A_k$ form a partition of $\Omega$, then
        $P(B) = \sum_{i=1}^{k} P(B \mid A_i)\,P(A_i)$.
  \item If $P(B \mid A) = P(B)$ (with $P(A), P(B) > 0$), then $P(A \cap B) = P(A)P(B)$ and $P(A \mid B) = P(A)$:
        independence is symmetric.
  \item The chain rule $P(A \cap B \cap C) = P(A)\,P(B \mid A)\,P(C \mid A \cap B)$. Use it to find the probability
        that three socks drawn from the drawer of Task 7(c) are all black.
  \item $P(A \mid B) > P(A)$ holds exactly when $P(B \mid A) > P(B)$: evidence that makes $A$ more likely is evidence
        that $A$ makes more likely.
  \item Simpson's paradox needs unequal groups: suppose that in every department the admission rate of women is at
        least that of men, and that women and men spread over the departments in the same proportions. Show with the
        law of total probability that the overall rate of women is then at least that of men.
  \item The odds form: show that $\frac{P(A \mid B)}{P(A^c \mid B)} = \frac{P(A)}{P(A^c)} \cdot
        \frac{P(B \mid A)}{P(B \mid A^c)}$ (posterior odds = prior odds $\times$ likelihood ratio). If two tests are
        independent given the person's true state, show that their likelihood ratios multiply, and use this to check
        Task 6(b)(iii).
\end{enumerate}

\medskip
\noindent Hand in your handwritten pages, with your name on each page.
